\chapter{命题与题目质量}

\section{让模型出题的风险}

题库只有四十七道题，每个知识点三道，难度档位有限。
学习者水平落在两道题之间时，只能取一道偏难或偏易的凑合，
而太难必错、太易必对，那一问的信息量就白费了。
因此引入模型现场命制变式题目，把题库当作种子而非天花板。

\begin{keynote}{先说风险，这条比什么都重要}
\emph{模型出题等于模型定义“正确答案”，这是本项目中后果最重的一类错误。}

讲义里写错一句，学习者可能看得出来，审核也拦得住。但答案标错了，
系统会\emph{拿一把错的尺子去量人}：学习者答对却被判错，掌握度往下走，
系统随后给他推送更简单的内容；学习者答错却被判对，盲区被漏掉，该补的没补。

错误不会停留在一句话上，而是\emph{灌进学情模型}，决定后面所有资源。
学习者还无法申诉，因为他并不知道标准答案是编出来的。

所以规矩是：\emph{生成题必须过闸，一道都不例外；过不了就丢弃，回退固定题库。
宁可用一道糙题，不能用一道错题。}
\end{keynote}

\section{五关审核}

命题审核的五道关卡见表 \ref{tab:vet}。前四关问的是“这道题对不对”，
第五关问的是“这道题好不好用”，后者详见 6.4 节。

\begin{table}[htbp]\centering\small
\caption{命题的五关审核}\label{tab:vet}
\begin{tabular}{@{}p{0.8cm}p{4.4cm}p{6.4cm}@{}}
\toprule
关 & 查什么 & 为什么 \\
\midrule
1 & 正确答案被所引切片支撑 & 沿用审核裁判那套判据，不给命题开小灶 \\
2 & 每个干扰项都\emph{不}成立 & 一题两解比答案错了还糟，且没有环节会发现 \\
3 & 题干不含切片以外的数值 & 只查题干，\emph{不查选项} \\
4 & 选项不重复、无长度线索 & 最长的恰好是答案，学习者不会也能蒙对 \\
5 & 结构质量 & 见 6.4 节 \\
\bottomrule
\end{tabular}
\end{table}

\begin{keynote}{踩坑：数值型选项不能用文本覆盖率判}
“0.5”这样的纯数字选项，按字符二元组切分后只剩“0”和“5”，
只要切片正文任意位置出现过这两个字符，覆盖率就是满分，判据完全失效。
数值型选项的正确判据是“这个数在不在知识库里”，黑白分明，
不该绕道文本相似度。因此审核对数值型与文字型选项分开处理。
\end{keynote}

第三关只查题干不查选项，这一点同样是踩出来的。
干扰项按定义就该带上知识库里没有的数值，
拿这一条去卡选项等于禁止出错误选项，第一版就是这么把所有题目全部毙掉的。
同理，长度线索规则跳过纯数值选项：一万比二长，
但没有哪个学习者会因此认为它是答案。

\section{否定感知：最隐蔽的一个坑}

在安全与规程类内容里，\emph{最好的干扰项恰恰是手册明令禁止的做法}。
知识库中记载“切勿在未按住解除按钮时强行驱动”，
那么“直接强行反向驱动”就是一个再标准不过的干扰项。
可它与该切片的字符二元组覆盖率高达 $0.60$，被判成“可能同样成立”而毙掉。

问题在于二元组重合看不见否定，“强行驱动”与“切勿强行驱动”在词面上几乎一样。
照这样卡下去，模型只能去造一眼就假的干扰项，题目的区分度随之归零。

修法是把切片\emph{按分句}切开，比较选项与“禁止句”和与普通正文的重合度，
更像禁止句就放行。这里用比较而不是绝对阈值，
因为干扰项通常只借用禁止句的一部分措辞，实测那一句只有 $0.40$，
阈值卡在 $0.5$ 会漏，卡在 $0.3$ 又会误判，比较式天然免疫这个标定问题。

两条反向保护同样重要。其一，选项\emph{自身}带有禁止词说明它在复述规则、本身成立，
不享受豁免。其二，切分粒度必须到分句：
“加注时必须打开排脂口，禁止在封闭状态下加注”这句话若只按句号切分，
整句都会被当成禁止句，连正确做法都被判成“被禁止”，
等于把正确答案塞进了干扰项。

\section{题目质量：题不好，尺子就不准}

命题闸门管的是“题\emph{对不对}”，题目质量模块管的是另一件事：\emph{题好不好用}。

两者不能互相替代。一道题可以完全正确、出处清楚、干扰项都不成立，
却仍然测不出任何东西：太简单人人都对，干扰项一眼就假，
两个干扰项几乎一样。这些不是错题，闸门放它们过去是对的。
但拿它们去估计掌握度，蒙对率就不再是四分之一，
估出来的掌握概率整体偏高，\emph{而且偏多少不知道}。

\subsection{上线第一天就抓到的两个问题}

\begin{keynote}{四十七道题里三十八道答案在 B}
学习者一路选 B 就能得到八成以上的分数，整套测评的测量效力等于零。
这是项目自己手写的题库，写的时候谁也没注意到。
现已用确定性置换打散，同时把三份画像里存储的作答下标一并重映射——
不同步重映射的话，对错模式会静默改变，
得到一批“通过但含义已变”的测试，是最难排查的一类问题。
\end{keynote}

\begin{keynote}{所有判断题的答案都是“正确”}
这一个更糟：位置偏斜至少还有四分之一的蒙对下限，答案全同则是百分之百。
根因是题干由已过审的断言直接生成，自然全是正命题。

修法是\emph{把审核闸反过来用}：扰动断言中的数值或特征术语，
再送回审核裁判，\emph{只有被判定为与知识库冲突才采纳}。
也就是说，这条假命题的“假”是知识库认证过的，不是我们自己认为它假。
认证不了就不产出反向题。现在真假比为八比四。
\end{keynote}

\subsection{结构检查与实测标定}

结构检查不需要作答数据，出题当场就能算，检查项见表 \ref{tab:flaw}。

\begin{table}[htbp]\centering\small
\caption{题目结构瑕疵检查项}\label{tab:flaw}
\begin{tabular}{@{}p{2.6cm}p{4cm}p{5cm}@{}}
\toprule
代码 & 查什么 & 为什么 \\
\midrule
\cd{LENCLUE} & 正确答案明显更长 & 应试技巧就能蒙对 \\
\cd{CATCHALL} & “以上都对”一类兜底项 & 测的不再是本知识点 \\
\cd{ABSOLUTE} & 干扰项绝对化措辞 & “所有”“绝不”会被直接排除 \\
\cd{DUPDISTRACT} & 两个干扰项几乎一样 & 名义四选一、实际三选一 \\
\cd{FARDISTRACT} & 干扰项离题 & 一眼能排除，白占一个选项位 \\
\cd{SHORTSTEM} & 题干过短 & 缺少作答语境 \\
\bottomrule
\end{tabular}
\end{table}

判定干扰项重复的阈值取 $0.60$ 而非直觉上的 $0.8$，这里有个标定细节：
中文短选项用字符二元组的 Jaccard 相似度时，改一个字就会掉两个二元组。
“由三个直线轴构成”与“由三个直线轴组成”实际是同一个选项，相似度却只有 $0.667$；
而真正不同的两个干扰项是 $0.304$。两档之间间隔足够大，$0.60$ 落在中间。

实测标定需要作答数据，包含两项指标。\emph{难度}用 $p$ 值表示，即答对比例，
注意方向与直觉相反——$p$ 越大题目越简单，报告中要写清楚，避免读反。
\emph{区分度}用点二列相关系数表示，含义是答对这道题的人在整份测评上是否也考得更好。

\begin{keynote}{区分度是发现答案标错的唯一可靠手段}
对本项目尤其要紧：生成题的答案是模型断言的，标错了不会有任何环节报错。
一旦标错，水平高的人反而更容易选到真正正确的那个选项而被判为错，
于是相关系数\emph{变成负数}。

判读标准为：大于 $0.30$ 区分度良好；$0.15$ 到 $0.30$ 可用；
$0$ 到 $0.15$ 几乎不提供信息；\emph{小于 $0$ 时应优先怀疑答案标错}，
而不是认为题目难。模拟验证显示，一道答案标错的题目区分度为 $-0.149$，
被直接判为不可用。
\end{keynote}

题库质量报告的当前结果为：共四十七道，平均质量分 85.7，判为不可用零道，
答案位置分布均匀，\emph{待改十九道}，主要问题是干扰项离题与长度线索。
\emph{这部分是人的活}——改写干扰项使其长度接近、具备干扰力，
属于命题技术，模型代劳不了。

\bibliographystyle{gbt7714-2005}
\bibliography{refs}
